\subsection{Pregunta experimental}

Con Feature Table v1, la capa agronómica y la capa empírica disponibles, Checkpoint~03C.1
separa el efecto de la \emph{representación} del efecto de la familia de modelo. Para hacerlo,
usa exactamente los mismos dos estimadores de referencia ---Ridge10 y ExtraTrees--- sobre
siete representaciones por track y conserva los folds de Checkpoint~02.

El diseño evita una comparación confusa del tipo representación A + modelo muy ajustado contra
representación B + modelo distinto. Aquí la pregunta es:
\[
\Delta_{\mathrm{representation}}
=
\RMSE(M,X^{(r_2)})-\RMSE(M,X^{(r_1)})
\]
para un mismo modelo $M$ y un mismo protocolo.

\subsection{Representaciones evaluadas}

\begin{table}[H]
\centering
\caption{Representaciones de 03C.1.}
\label{tab:cp03c1-reps}
\small
\begin{tabular}{llrr}
\toprule
Track & Representación & $p$ & Discovery \\
\midrule
clean & B0 base & 692 & no \\
clean & B1 agronomic & 123 & no \\
clean & B2 empirical & 91 & no \\
clean & B3 base+agro & 815 & no \\
clean & B4 base+empirical & 783 & no \\
clean & B5 all & 906 & no \\
clean & B7 all+discovery & 906 & sí, hasta 20 \\
\midrule
competition & B0 base & 1,400 & no \\
competition & B1 agronomic & 350 & no \\
competition & B2 empirical & 335 & no \\
competition & B3 base+agro & 1,750 & no \\
competition & B4 base+empirical & 1,735 & no \\
competition & B5 all & 2,085 & no \\
competition & B7 all+discovery & 2,085 & sí, hasta 20 \\
\bottomrule
\end{tabular}
\end{table}

En B7 las columnas base son fijas, pero las expresiones supervisadas no lo son: cada outer
training fold descubre sus propias expresiones. El número 20 indica un máximo por fit, no un
set global de 20 fórmulas preseleccionadas.

\subsection{Métrica OOF versus promedio de folds}

03C.1 conserva dos resúmenes que no son idénticos cuando los folds tienen tamaños distintos.
Si $F$ es el conjunto de folds:
\[
\overline{\RMSE}_{\mathrm{fold}}
=
\frac{1}{|F|}
\sum_{f\in F}\RMSE_f,
\]
mientras
\[
\RMSE_{\mathrm{OOF}}
=
\sqrt{
\frac{1}{138}
\sum_{i=1}^{138}
(y_i-\hat y_i^{\mathrm{OOF}})^2
}.
\]
El primero pondera cada fold por igual; el segundo pondera cada parcela por igual. En
municipality-grouped, donde los folds pueden ser muy desiguales, reportar ambos evita esconder
la influencia de un municipio grande.

\subsection{Resultados en state-stratified}

Para competition + ExtraTrees, B7 obtiene OOF RMSE 0.4992, B0 0.5003, B1 0.5041 y B5
0.5045. Es decir, bajo este protocolo las diferencias entre representaciones ricas son
pequeñas; la capa empírica sola B2 es claramente peor, con 0.6470.

Con Ridge10 el comportamiento cambia: B0 alcanza 0.6534, B4 0.6587, B3 0.6644, B7 0.6679 y
B1 0.6809. La compactación agronómica que puede ser razonable para árboles no mejora
automáticamente un Ridge de alpha fijo en este track.

En clean, ExtraTrees B7 alcanza 0.4948 OOF y B3 0.5021, mientras Ridge sobre la capa
agronómica compacta B1 alcanza 0.5888 frente a 0.6718 de B0. Esto muestra que el valor de una
representación depende fuertemente de la inductive bias del modelo.

\subsection{Resultados en municipality-grouped}

Bajo el protocolo espacialmente más exigente, competition + ExtraTrees favorece B1 agronomic:
\[
\RMSE_{\mathrm{OOF}}=0.8590,
\]
seguido por B7=0.8675 y B0=0.8974. Para Ridge, B0 es claramente mejor:
0.8059, frente a B4=0.8464, B3=0.8676, B5=0.8976 y B7=0.9142.

En clean, ExtraTrees B7 llega a 0.8488 y B1 a 0.8762; Ridge B0 queda en 0.8741. Ninguna
representación produce una mejora universal. La diferencia entre protocolos sigue siendo más
grande que muchas diferencias internas de representación.

\subsection{Lección de 03C.1}

Tres conclusiones pasan a 03C.2:

\begin{enumerate}
\item \textbf{C0/base debe sobrevivir}: pese a la alta dimensionalidad, sigue siendo una
representación competitiva para modelos lineales latentes como PLS/Ridge.
\item \textbf{C1/agronomic merece un benchmark propio}: su compactación funciona bien con
árboles bajo el stress grouped y ofrece una representación de sólo 350 columnas.
\item \textbf{El bulk empírico no se añade por default}: B2 solo es débil y B5 puede
degradar algunos modelos; las primitivas empíricas más útiles se retienen de forma controlada
para discovery.
\end{enumerate}

Por ello, 03C.2 reduce el espacio de representaciones en vez de continuar acumulando capas.
Este es un patrón deliberado del proyecto: cada checkpoint debe \emph{cerrar} grados de
libertad antes de abrir una búsqueda más costosa.
